% The set-up of a search: the checks of the input and the refusals, the squarefree binary form, the
% reversal and its map of the points, the remainder chain, the points at infinity, and the exit
% when the curve has no points modulo a prime.
\chapter{The set-up}\label{chap:setup}

This chapter describes what a search does with its input before the sieve: it checks the coefficient
list and the height bound, drops the leading zero coefficients, refuses a binary form that is not
squarefree, may reverse the polynomial, delivers the points at infinity, and ends when the curve has
no points modulo one of the primes it examines. The refusal theorems (Theorem~\ref{thm:refusals})
rest on it, and so do the statements on the points delivered, since after a reversal \ratpoints{}
searches another polynomial and maps its points back. The rest of the Sturm step, which begins with
the remainder chain of this chapter, is in Chapter~\ref{chap:region}, the tests on the denominators
that follow the set-up in Chapter~\ref{chap:denominators}, and the loop over the primes in which the
exit fires in Chapter~\ref{chap:moduli}.

Throughout, $c = [c_0, \dots, c_n]$ is the coefficient list, constant term first; $d$ is its degree,
the largest index with $c_d \ne 0$; $D = 2 \lceil d/2 \rceil$ and $N = 2 \lceil n/2 \rceil$ are~$d$
and~$n$ rounded up to even numbers; $F_c(x, z) = \sum_{i \le d} c_i x^i z^{D-i}$ is the binary form
of the points (Definition~\ref{def:binaryForm}), $\Phi_c = \sum_{i \le n} c_i X^i Z^{N-i}$ the
binary form of the refusals (Definition~\ref{def:binaryFormPoly}), and $f = \sum_i c_i X^i \in
\QQ[X]$ the polynomial that the C code works with.

\section{The order of the checks}\label{sec:setupOrder}

The program refuses more than $\mathrm{maxDegree} + 1$ coefficients with its error~3, fewer than two
with error~5, and a height bound below~$1$ or one that is not a \verb|long| with error~4, before it
calls the library (\csource{main.c}{362}{380}). The library refuses a degree above
$\mathrm{maxDegree}$ (\csource{find_points.c}{4557}{4558}), and then:
\begin{enumerate}
  \item it drops the leading zero coefficients and refuses the form as not squarefree when that
    lowers $\lceil n/2 \rceil$, that is when $d + 2 \le N$, so that $Z^2 \mid \Phi_c$
    (\csource{find_points.c}{4624}{4632}); this catches every list of two or more entries of
    degree~$0$, so that the refusal of a degree~$0$ as bad arguments that follows
    (\csource{find_points.c}{4633}{4633}) is never met from the program; then it refuses a height
    bound below~$1$ (\csource{find_points.c}{4658}{4658});
  \item it turns the reversal off when the caller gives search intervals or a range of the
    denominators other than $1, \dots, H$ (\csource{find_points.c}{4686}{4687}), and runs the
    reversal check, which refuses the form when~$d$ is even and $c_0 = c_1 = 0$, so that
    $X^2 \mid \Phi_c$, and may reverse the polynomial (\csource{find_points.c}{4732}{4803});
  \item it takes the default interval $[-H, H]$, its ends rounded outward to doubles, when the
    caller gives none, and otherwise refuses as bad arguments search intervals without $l \le u$,
    not increasing and strictly apart, or with a NaN end (\csource{find_points.c}{4831}{4855}),
    which no command line causes;
  \item with the Sturm step on, it refuses the form when the remainder chain of the polynomial
    searched meets a zero remainder, then restricts the search intervals, and returns when nothing
    is left to search and there is no point at infinity (\csource{find_points.c}{4857}{4892});
  \item it sets up the tests on the denominators and chooses the moduli, each with an exit when it
    finds that the curve has no points (\csource{find_points.c}{4954}{4963},
    \csource{find_points.c}{4996}{5012});
  \item it delivers the points at infinity (\csource{find_points.c}{5068}{5105}) and sieves.
\end{enumerate}
The program reports the refusal as not squarefree with its error~8, and the other negative codes of
the library with its error~9, but for a lack of memory, error~10 (\csource{main.c}{256}{261}). All
refusals come before the first call of the callback. Where the program and the library differ, the
model follows the program: it refuses the coefficient list and the height bound before the leading
zeros (so that \verb|5 0| with $H = 0$ is error~4, as in the program), and a height bound
of~$2^{63}$ or more, which the library cannot be given, with error~4. These first checks are in the
search with a callback (Definition~\ref{def:findPoints}); the rest of the set-up is one definition.
Appendix~\ref{chap:correspondence} lists how else the components differ from the C code; the main
deviations of the set-up are noted in the sections below.

\begin{definition}[The set-up of a search]\label{def:setupSearch}
  \uses{def:Options, def:reversal, def:maxDegree, def:deliver, def:callback, def:computeSturm, def:libraryReport, def:searchIntervals, def:searchPoints, def:tableOfPrime}
  \lean{Ratpoints.findPointsWork}
  \leanok
  The refusals (Definition~\ref{def:Options}) carry the errors~3, 5, 4, 8 and~9 of the program; an
  integer is a \verb|long| when it lies in $[-2^{63}, 2^{63})$. The set-up of \verb|find_points_work_1|, whose state is declared at
  \csource{find_points.c}{4577}{4585}, for the list~$c$ of degree~$d$ and the options~$o$,
  normalizes the options (\lname{Ratpoints.Options.normalize}), so that the reversal is allowed only
  when~$o$ allows it and its search region is the whole box; runs the reversal check
  (Definition~\ref{def:reversal}), which may refuse or hand on a list~$c'$ of degree~$d'$; takes
  the search intervals of~$o$, or the default one, refusing malformed ones
  (\lname{Ratpoints.searchIntervals}); runs the Sturm step on~$c'$ (\lname{Ratpoints.sturmRegion}),
  which with the step on refuses the form when the remainder chain of~$c'$ is none, returns at once
  with nothing delivered when it keeps no interval and~$c'$ has no point at infinity, and otherwise
  narrows the intervals; and runs the search of~$c'$ (\lname{Ratpoints.searchPoints}) with the flag
  of the reversal and the options of the output as the flags of the check. Whether the coefficients
  of~$c'$ are all \verb|long|s, which the C code tests in the set-up and the model tests in the
  search of~$c'$, where it is used (Definition~\ref{def:searchPoints}), chooses how the tables
  reduce them modulo a prime: by the helper \verb|mod| of the C code
  (\lname{Ratpoints.Sieve.modLong}) or by the Euclidean remainder, which agree
  (Chapter~\ref{chap:rows}).
\end{definition}

\begin{lemma}[The set-up in the order of the program]\label{lem:setupOrder}
  \uses{def:findPoints, def:setupSearch, def:reversal, def:maxDegree, def:regionSpec, def:computeSturm, def:searchIntervals}
  \lean{Ratpoints.findPointsWith_eq, Ratpoints.findPointsWith_eq_ok_iff, Ratpoints.findPointsWith_eq_error_iff, Ratpoints.fitsLong_natCast, Ratpoints.IsSearchedPoly}
  \leanok
  With the print channel that prints nothing, the search with a callback refuses, in this order,
  more than $\mathrm{maxDegree} + 1$ coefficients, fewer than two, $H = 0$ or $H \ge 2^{63}$,
  $\lceil d/2 \rceil < \lceil n/2 \rceil$ (error~8), the refusal of the reversal check with the
  switch of the normalized options, malformed search intervals, and the refusal of the Sturm step;
  otherwise it returns the outcome of the search of the list~$c'$ that the check hands on, or, when
  the Sturm step returns before the search, the data of the callback with the count~$0$, the flag of
  the reversal and no interval. The refusals do not depend on the callback. A natural number is a
  \verb|long| exactly when it is below~$2^{63}$. The predicate $\mathrm{IsSearchedPoly}$ collects
  what the refusals leave for a successful search, from which the main theorems start
  (Chapter~\ref{chap:theorems}): the reversal check of the normalized options hands on~$c'$ with its
  degree~$d'$; $1 \le d' \le \mathrm{maxDegree}$; $1 \le H < 2^{63}$; the caller's search
  intervals are well formed (Definition~\ref{def:regionSpec}); and with the Sturm step on, the
  remainder chain of~$c'$ is not none.
\end{lemma}

\begin{proof}
  \leanok
  Unfolding the definitions gives the order; a print channel that prints nothing leaves the data as
  they are, and every refusing branch returns the same error for any callback. A natural number is
  never below~$-2^{63}$. After the refusals~$c$ has at least two entries and the leading zeros did
  not lower $\lceil n/2 \rceil$, so $d \ge 1$; the reversal check returns the degree of its list,
  never raises it and keeps it at least~$1$ (Lemma~\ref{lem:reversal_swap}).
\end{proof}

\section{The squarefree binary form}\label{sec:setupSquarefree}

As $\Phi_c = Z^{N-d}\,\Phi^0$, where~$\Phi^0$ is the homogenization of~$f$ in degree~$d$, which~$Z$
does not divide, $\Phi_c$ is squarefree exactly when $N \le d + 1$ and~$f$ is squarefree
(Lemma~\ref{lem:squarefree_binaryFormPoly}): the first condition is the test of the leading zeros,
the second is what the remainder chain tests (Section~\ref{sec:remainderChain}). Mathlib has nothing
on the squarefreeness of a homogenization; the facts that are needed are proved at Mathlib's
generality in \verb|Ratpoints/Mathlib|. The search itself needs no squarefreeness: the statements on
the points hold for every polynomial that the set-up hands on, also for a form that is not
squarefree, which the search accepts when the Sturm step is off.

\begin{lemma}[Squarefreeness in the terms of the C code]\label{lem:squarefree_binaryFormPoly}
  \uses{def:binaryForm, def:binaryFormPoly}
  \lean{Polynomial.squarefree_homogenize_iff, Polynomial.homogenize_reflect, Polynomial.squarefree_homogenize_iff_squarefree_reflect, Ratpoints.poly, Ratpoints.natDegree_poly, Ratpoints.binaryFormPoly_eq_homogenize, Ratpoints.squarefree_binaryFormPoly_iff, Ratpoints.not_squarefree_binaryFormPoly_of_lt, Ratpoints.not_squarefree_binaryFormPoly_of_getD_zero}
  \leanok
  \begin{enumerate}
    \item Over a domain and for $m \ge \deg p$, the homogenization of degree~$m$ of a
      polynomial~$p$ is squarefree if and only if~$p$ is squarefree and $m \le \deg p + 1$. The
      homogenization of degree~$m$ of the reflection of~$p$ in degree~$m$ is that of~$p$ with the
      variables exchanged; so, for $m \ge \deg p$, the homogenization of degree~$m$ of~$p$ is
      squarefree exactly when the reflection of~$p$ in degree~$m$ is, provided that~$m$ exceeds
      the degree of the reflection by at most one.
    \item The polynomial~$f$ of the list~$c$ has the degree~$d$ of~$c$, and~$\Phi_c$ is its
      homogenization of degree~$N$; so~$\Phi_c$ is squarefree if and only if~$f$ is squarefree and
      $N \le d + 1$. In particular~$\Phi_c$ is not squarefree when the leading zeros lower
      $\lceil n/2 \rceil$, nor when $c_0 = c_1 = 0$, whatever the degree.
  \end{enumerate}
\end{lemma}

\begin{proof}
  \leanok
  (1) The homogenization of degree~$m$ is that of degree~$\deg p$ times $Z^{m - \deg p}$. If
  $m \ge \deg p + 2$, then~$Z^2$ divides it, and if $a^2 \mid p$, the square of the homogenization
  of~$a$ divides it. Conversely, if~$A^2$ divides it and $m \le \deg p + 1$, then~$A$ is
  homogeneous, as a factor of a homogeneous polynomial over a semiring without zero divisors
  (\lname{MvPolynomial.IsHomogeneous.of\_mul\_left}), $Z \nmid A$, and the square of $A(X, 1)$
  divides~$p$, so $A(X, 1)$ is a constant, and then so is~$A$. The reflection compares
  coefficients, and the exchange of the variables is a ring automorphism, which keeps
  squarefreeness. (2) The coefficient of~$X^i$ in~$f$ is~$c_i$, so~$f$ has the degree $d \le N$;
  $\Phi_c$ and the homogenization agree term by term, and~(1) applies. A drop of
  $\lceil n/2 \rceil$ means $d + 2 \le N$, and $c_0 = c_1 = 0$ gives $X^2 \mid f$.
\end{proof}

\section{The reversal}\label{sec:reversal}

When the caller allows it and the search region is the whole box, \ratpoints{} may replace the
polynomial by its reversal, in three cases in which it expects the search of the reversed polynomial
to be faster. The reversal exchanges~$x$ and~$z$ in the binary form; the C code maps each point of
the reversed curve back to the caller's curve before it delivers it, so that the output holds the
same points in another order, which the specification of the order follows: it names the polynomial
searched, as the reversal check chooses it (Definition~\ref{def:orderedPoints}). The model makes a
new list where the C code reverses its array in place, and writes back to the caller the flag of the
reversal, not the reversed coefficients.

\begin{definition}[The reversal check and the map of its points]\label{def:reversal}
  \uses{def:Options}
  \lean{Ratpoints.Setup.reversal, Ratpoints.Setup.isSquare, Ratpoints.Sieve.swapPoint}
  \leanok
  For the switch of the reversal and the list~$c$ of degree~$d$: if the switch is off, $c$ itself;
  otherwise, for even~$d$ with $c_0 = 0$, the refusal (not squarefree) if also $c_1 = 0$, and else
  the list $c_d, \dots, c_1, c_0$ with the degree~$d - 1$ (case~1); for even~$d$ with~$c_d$ a square
  and~$c_0$ not, the list $c_d, \dots, c_0$ (case~2); for odd~$d$ with $|c_d| \ne 1$, $c_0 = 0$ and
  $|c_1| = 1$, the list $c_0, c_d, \dots, c_1$ (case~3); and otherwise~$c$ itself; with a flag that
  says whether the list is reversed (\csource{find_points.c}{4744}{4756},
  \csource{find_points.c}{4759}{4770}, \csource{find_points.c}{4775}{4790}). The squares are those
  of \verb|mpz_perfect_square_p|: $0$ and the squares of the positive integers. The map of a
  reversed search sends a point $(a, y, b)$ of the reversed curve to $(b, y, a)$ for $a \ge 0$ and
  to $(-b, y, -a)$ for $a < 0$; the exact check applies it to every point it finds when the flag of
  the reversal is set (Chapter~\ref{chap:sieve}), and the points at infinity use its image of the
  candidate $(1, 0)$ (Definition~\ref{def:pointsAtInfinity}).
\end{definition}

\begin{lemma}[The reversal exchanges~$x$ and~$z$]\label{lem:reversal_swap}
  \uses{def:reversal, def:binaryForm, def:binaryFormPoly, def:regionSpec, def:Options, def:searchIntervals, lem:squarefree_binaryFormPoly}
  \lean{Ratpoints.Setup.exists_reversal_eq_error_iff, Ratpoints.exists_reversal_normalize_eq_error_iff, Ratpoints.isFullBox_of_reversal_normalize_eq_ok, Ratpoints.Setup.degree_eq_of_reversal_eq_ok, Ratpoints.Setup.degree_le_of_reversal_eq_ok, Ratpoints.Setup.one_le_of_reversal_eq_ok, Ratpoints.Setup.getD_eq_of_reversal_eq_ok, Ratpoints.Setup.binaryForm_eq_of_reversal_eq_ok, Ratpoints.Setup.squarefree_binaryFormPoly_iff_of_reversal_eq_ok}
  \leanok
  The check refuses exactly when its switch is on, $d$ is even and $c_0 = c_1 = 0$; with the switch
  of the normalized options, exactly when the options allow the reversal, their search region is the
  whole box (Definition~\ref{def:regionSpec}), $d$ is even and $c_0 = c_1 = 0$; and it reverses only
  for the whole box. For~$d$ the degree of~$c$: the degree returned is that of the list returned, at
  most~$d$, and at least~$1$ when $d \ge 1$. If the check reverses, the list~$c'$ it returns is the
  reflection of~$c$ in degree~$D$, $c'_i = c_{D-i}$ for $i \le D$ and $c'_i = 0$ beyond, and
  $F_{c'}(x, z) = F_c(z, x)$ for all integers~$x$ and~$z$. If the leading zeros of~$c$ do not lower
  $\lceil n/2 \rceil$, then~$\Phi_c$ is squarefree if and only if the polynomial of the list handed
  on is, reversed or not.
\end{lemma}

\begin{proof}
  \uses{lem:squarefree_binaryFormPoly, lem:searchIntervals}
  \leanok
  By the cases of the check. In each case of a reversal the list built is the reflection, with
  $c_{d+1} = 0$ in case~3; the new leading coefficient, $c_1$, $c_0$ or~$c_1$, is not zero, so the
  reversed list has the degree returned and the same~$D$, and
  $F_{c'}(x, z) = \sum_i c_{D-i}\, x^i z^{D-i} = F_c(z, x)$. Without a reversal
  Lemma~\ref{lem:squarefree_binaryFormPoly} gives the last claim; after one, the polynomial of~$c'$
  is the reflection of~$f$ in degree $D = N$, of degree~$D - 1$ or~$D$, and the reflection does not
  change whether the homogenization of degree~$D$ is squarefree
  (Lemma~\ref{lem:squarefree_binaryFormPoly}).
\end{proof}

\begin{lemma}[The points of a reversed search]\label{lem:swapPoint_bijection}
  \uses{def:reversal, def:IsPoint, def:bruteForce, thm:bruteForce, def:orderedPoints, def:maxDegree}
  \lean{Ratpoints.IsNormalizedPair.swapPoint, Ratpoints.Sieve.swapPoint_inj, Ratpoints.Sieve.isPoint_swapPoint, Ratpoints.Sieve.perm_map_swapPoint_bruteForce, Ratpoints.searchedPoly_eq_of_reversal_eq_ok, Ratpoints.fitsLong_neg}
  \leanok
  The map of a reversed search keeps the normalized pairs $(x, z)$ of height at most~$H$
  (Definition~\ref{def:IsPoint}) and is injective on the triples with $z > 0$ or $(x, z) = (1, 0)$.
  For lists~$c$ and~$c'$ with $F_{c'}(x, z) = F_c(z, x)$ for all~$x$ and~$z$, it sends the
  normalized points of height at most~$H$ of~$c'$ to points of~$c$, and the brute-force reference
  of~$c'$, mapped, is a permutation of that of~$c$. The reversal check hands on the polynomial and
  the flag of the specification of the order (Definition~\ref{def:orderedPoints}). The negation of
  an integer of absolute value below~$2^{63}$ is a \verb|long|, so the map stays among the
  \verb|long|s for $H < 2^{63}$.
\end{lemma}

\begin{proof}
  \leanok
  $F_c(-b, -a) = F_c(b, a)$, since~$D$ is even. For $a \ne 0$ the image has $z = |a| > 0$ and
  $|x| = b$; for $a = 0$ the normalization forces $b = 1$, and the image $(1, y, 0)$ is a point at
  infinity; so the box goes onto the box. On the triples with $z > 0$ or $(x, z) = (1, 0)$ the map
  is its own inverse; both references list their points once (Theorem~\ref{thm:bruteForce}), so the
  mapped list is a permutation.
\end{proof}

\section{The remainder chain}\label{sec:remainderChain}

\begin{definition}[The remainder chain]\label{def:remainderChain}
  \lean{Ratpoints.Sturm.chain, Ratpoints.Sturm.IsChain, Ratpoints.Sturm.IsChainResult, Ratpoints.Sturm.chain_spec, Ratpoints.Sturm.IsChain.size_le}
  \leanok
  For the coefficients $c_0, \dots, c_d$ of~$f$, with $d \ge 1$ and $c_d \ne 0$, the rows of
  integers that \verb|sturm.c| computes (\csource{sturm.c}{99}{163}): $s_0 = f$, $s_1 = f'$, and for
  $k = 2, \dots, d$ the row~$s_k$ from the remainder~$t$ of~$s_{k-2}$ by~$s_{k-1}$: starting from
  $t = s_{k-2}$, while $\deg t \ge \deg s_{k-1}$, $t$ becomes
  $u t - v X^{\deg t - \deg s_{k-1}} s_{k-1}$, with~$u$ and~$v$ the leading coefficients
  of~$s_{k-1}$ and~$t$ divided by their~$\gcd$ and signed so that $u > 0$, and its zero leading
  coefficients are dropped; the result is none when~$t$ becomes zero. Otherwise~$s_k$ is~$-t$
  divided by the~$\gcd$ of its coefficients, and the chain ends when~$s_k$ is constant. The
  contract, proved with the program logic of Lean core: the loops end, and the result is a chain
  ($\mathrm{IsChain}$: rows as just described, each with a nonzero leading coefficient, their sizes
  decreasing from~$s_1$ on) whose last row is constant, or none, and then some chain has a last row
  of positive degree that reduces the row before it to zero. A chain has at most~$d + 1$ rows, the
  size of the table of the C code.
\end{definition}

\begin{lemma}[The chain detects a repeated factor]\label{lem:chain_squarefree}
  \uses{def:remainderChain, lem:squarefree_binaryFormPoly, def:reversal}
  \lean{Ratpoints.Sturm.IsChain.exists_mul_eq, isCoprime_iff_of_mul_eq_mul_sub_mul, Ratpoints.Sturm.run_chain_eq_none_iff, Ratpoints.Sturm.run_chain_eq_none_iff_of_reversal_eq_ok}
  \leanok
  Three consecutive rows of a chain satisfy $A\, s_k = Q\, s_{k+1} - g\, s_{k+2}$ with integers
  $A, g > 0$ and some $Q \in \QQ[X]$. For $d \ge 1$ and $c_d \ne 0$, the chain is none if and only
  if~$f$ is not squarefree over~$\QQ$. So for a list~$c$ of at least two entries whose leading zeros
  do not lower $\lceil n/2 \rceil$, the chain of the list that the reversal check hands on, with the
  degree it returns, is none if and only if~$\Phi_c$ is not squarefree.
\end{lemma}

\begin{proof}
  \uses{lem:reversal_swap}
  \leanok
  The relation is read off the arrays coefficient by coefficient. In a commutative ring, for
  units~$a$ and~$b$, $a x = q y - b z$ makes~$y$ and~$z$ coprime exactly when~$x$ and~$y$ are; so
  coprimality passes down the chain from $(f, f')$ to its last two rows, which a constant last row
  makes coprime and a last row of positive degree that reduces the row before it to zero makes not
  coprime. Over~$\QQ$, $f$ and~$f'$ are coprime exactly when~$f$ is separable, that is squarefree.
  The last claim follows with Lemma~\ref{lem:reversal_swap}: the degree handed on is at least~$1$,
  and~$\Phi_c$ is squarefree exactly when the polynomial handed on is.
\end{proof}

The C code keeps the rows, not only the verdict, because the Sturm step goes on with them: the
relation with positive~$A$ and~$g$ and the decreasing degrees make the rows positive multiples of
the Sturm sequence of~$f$ and~$f'$, whose sign changes the step counts (Chapter~\ref{chap:region}).
The model keeps each row as the array of its coefficients up to its degree, where the C code keeps a
table of GMP integers with an array of degrees.

\section{The points at infinity}\label{sec:pointsAtInfinity}

The points at infinity of the polynomial searched are the triples $(1, y, 0)$ with $y^2 = c_D$: for
odd~$d$ the one point $(1 : 0 : 0)$, for even~$d$ the points $(1 : \pm\sqrt{c_d} : 0)$ when~$c_d$ is
a square, and none otherwise. After a reversal they are the points $(0 : y : 1)$ of the caller's
curve.

\begin{definition}[The points at infinity]\label{def:pointsAtInfinity}
  \uses{def:deliver, def:reversal, def:bruteForce, def:callback, lem:deliverList}
  \lean{Ratpoints.Sieve.pointsAtInfinity, Ratpoints.Sieve.pointsAtInfinityList, Ratpoints.Sieve.pointsAtInfinity_eq_deliver, Ratpoints.Sieve.pointsAtInfinityList_eq_filter_map_pointsOfDen}
  \leanok
  For the callback, the list~$c$ of degree~$d$, the flags of the check and a delivery whose flag of
  the stop is clear: when~$d$ is odd or~$c_d$ a square, with the candidate $(a, b) = (1, 0)$, or
  $(0, 1)$ when the flag of the reversal is set, one call with $(a : 0 : b)$ under
  $\mathrm{noCheck}$ or for odd~$d$, and otherwise the call with $(a : s : b)$ for
  $s = \lfloor \sqrt{c_d} \rfloor$ and then, unless that call stopped the search or~$\mathrm{noY}$
  is set, the call with $(a : -s : b)$; nothing otherwise (\csource{find_points.c}{4894}{4896},
  \csource{find_points.c}{5068}{5074}, \csource{find_points.c}{5078}{5099}). The contract: on a
  delivery whose flag is clear these calls are the delivery (Definition~\ref{def:deliver}) of the
  list of $(1, 0, 0)$, or of $(1, s, 0)$ and, unless~$\mathrm{noY}$, $(1, -s, 0)$, each mapped by
  the map of the reversal when its flag is set, and of the empty list when there are no points at
  infinity; without $\mathrm{noCheck}$ and for odd~$d$ or $c_d \ne 0$, this list is that of the
  brute-force reference over the candidate $(1, 0)$, mapped after a reversal and without the point
  with $y < 0$ under~$\mathrm{noY}$.
\end{definition}

Like the C code, the model has no test $s \ne 0$, which would matter only for even~$d$ and $c_d = 0$
(the point $(1, 0, 0)$ twice), excluded by the refusals. The C code computes the flag \verb|lcfsq|
before the reversal check and clears it in case~2; the model tests~$c_d$ on the polynomial searched,
which agrees for even degree, the only one that reads the flag. The search delivers the points at
infinity after the choice of the moduli and returns when a call stops it there
(Definition~\ref{def:searchPoints}).

\section{The exit by a prime}\label{sec:exitByPrime}

While it chooses the moduli (Chapter~\ref{chap:moduli}), \ratpoints{} computes for each prime~$p$ it
examines the table of the residues~$x$ with~$f(x)$ a square modulo~$p$, with a last entry for the
points at infinity modulo~$p$ (Chapter~\ref{chap:rows}). When no entry is true,
\verb|examine_prime| returns~$-1$ (\csource{find_points.c}{1762}{1762}) and the search ends with no
points (\csource{find_points.c}{2613}{2614}, \csource{find_points.c}{2994}{2998},
\csource{find_points.c}{4996}{5012}). The model's table with the exit,
\lname{Ratpoints.Sieve.examinePrime?}, is then none, and otherwise the reference table
\lname{Ratpoints.Sieve.examinePrime}, for every prime below
$\mathrm{maxPrimeEven} = 2^{\mathrm{maxBitsInPrime}}$, the bound of the moduli of the sieve, below
which lie all the primes that the choice examines.

\begin{lemma}[A curve without points modulo~$p$ has none]\label{lem:exitByPrime}
  \uses{def:IsPoint, def:binaryForm, def:tableOfPrime, lem:hornerBarrett}
  \lean{Ratpoints.Sieve.any_isFSquare_examinePrime_of_sq_eq, Ratpoints.Sieve.isSquare_getD_degree_of_sq_eq_binaryForm, Ratpoints.Sieve.not_isPoint_of_run_examinePrime?_eq_none, Ratpoints.Sieve.not_isPoint_of_lt_numPrimes_of_run_examinePrime?_eq_none}
  \leanok
  Let~$p$ be a prime and~$d$ the degree of~$c$. Integers~$a$ and~$y$ and a natural number~$b$ with
  $\gcd(a, b) = 1$ and $y^2 = F_c(a, b)$ make some entry of the table of~$p$ true: for $p \nmid b$
  the entry of the residue~$a b^{-1}$, and for $p \mid b$ the entry at infinity, since then~$c_d$ is
  a square modulo~$p$ when~$d$ is even. So if $p < \mathrm{maxPrimeEven}$ and the table of~$p$ with
  the exit is none, $c$ has no normalized point of any height; this holds in particular for the
  prime at each position below $\mathrm{numPrimes}$ of the table.
\end{lemma}

\begin{proof}
  \uses{lem:sieveEntry_sound, lem:hornerBarrett}
  \leanok
  For $p \nmid b$ the row of~$b$ admits~$a$ (Lemma~\ref{lem:sieveEntry_sound}), so the entry
  of~$a b^{-1}$ is true. For $p \mid b$, $b = 0$ included, $p \nmid a$ and
  $y^2 \equiv c_D a^D \bmod p$, so the entry at infinity is true: for even~$d$ because~$c_d$ is then
  a square modulo~$p$, for odd~$d$ always. For $0 < p < \mathrm{maxPrimeEven}$ the table with the
  exit is none exactly when no entry is true (Lemma~\ref{lem:hornerBarrett}).
\end{proof}

When the exit fires the curve has no points at all, so it changes no output. After a reversal it
never fires: each case leaves the caller's curve a point of height~$1$, $(0 : 0 : 1)$ in cases~1
and~3 and the points at infinity in case~2.
